\begin{thebibliography}{99}
\bibitem[Araka et al.(2020)]{education_araka2020} Eric Araka; E. Maina; Rhoda Gitonga; Robert O. Oboko. Research trends in measurement and intervention tools for self-regulated learning for e-learning environments—systematic review (2008–2018). Research and Practice in Technology Enhanced Learning, 2020. DOI: 10.1186/s41039-020-00129-5. URL: \url{https://doi.org/10.1186/s41039-020-00129-5}. Consulta: 5 out. 2026. Nível de acesso registrado na matriz de fontes.


\bibitem[Baars e Viberg(2022)]{education_baars2022} Martine Baars; Olga Viberg. Mobile Learning to Support Self-Regulated Learning: A Theoretical Review. Int. J. Mob. Blended Learn., 2022. DOI: 10.4018/IJMBL.315628. URL: \url{https://doi.org/10.4018/IJMBL.315628}. Consulta: 5 out. 2026. Nível de acesso registrado na matriz de fontes.


\bibitem[Brahma e Saikia(2023)]{education_brahma2023} B. Brahma; P. Saikia. Influence of self-regulated learning on the academic procrastination of college students. Journal of Education and Health Promotion, 2023. DOI: 10.4103/jehp.jehp\_1106\_22. URL: \url{https://doi.org/10.4103/jehp.jehp_1106_22}. Consulta: 5 out. 2026. Nível de acesso registrado na matriz de fontes.


\bibitem[Broadbent et al.(2020)]{education_broadbent2020} Jaclyn Broadbent; E. Panadero; M. Fuller-Tyszkiewicz. Effects of mobile-app learning diaries vs online training on specific self-regulated learning components. Educational Technology Research and Development, 2020. DOI: 10.1007/s11423-020-09781-6. URL: \url{https://doi.org/10.1007/s11423-020-09781-6}. Consulta: 5 out. 2026. Nível de acesso registrado na matriz de fontes.


\bibitem[Chen et al.(2024)]{education_chen2024} Yi-Chun Chen; Gwo-Jen Hwang; Chiu-Lin Lai. Motivating students to become self-regulatory learners: A gamified mobile self-regulated learning approach. Education and Information Technologies, 2024. DOI: 10.1007/s10639-024-12462-z. URL: \url{https://doi.org/10.1007/s10639-024-12462-z}. Consulta: 5 out. 2026. Nível de acesso registrado na matriz de fontes.


\bibitem[Forstervold et al.(2022)]{education_forstervold2022} Knut Inge Forstervold; S. Ludvigsen; Helge I. Strømsø. Students’ time management and procrastination in the wake of the pandemic. Educational Psychology, 2022. DOI: 10.1080/01443410.2022.2102582. URL: \url{https://doi.org/10.1080/01443410.2022.2102582}. Consulta: 5 out. 2026. Nível de acesso registrado na matriz de fontes.


\bibitem[Han et al.(2023)]{education_han2023} Jaeyun Han; D. DiGiacomo; Ellen L. Usher. College students’ self-regulation in asynchronous online courses during COVID-19. Studies in Higher Education, 2023. DOI: 10.1080/03075079.2023.2201608. URL: \url{https://doi.org/10.1080/03075079.2023.2201608}. Consulta: 5 out. 2026. Nível de acesso registrado na matriz de fontes.


\bibitem[Hartley et al.(2022)]{education_hartley2022} Kendall Hartley; Emily Shreve; Dan Gianoutsos; Lisa D. Bendixen. Smartphone as a Self-regulatory Planning Tool: Promise or Peril. International journal of interactive mobile technologies, 2022. DOI: 10.3991/ijim.v16i14.28783. URL: \url{https://doi.org/10.3991/ijim.v16i14.28783}. Consulta: 5 out. 2026. Nível de acesso registrado na matriz de fontes.


\bibitem[Lim et al.(2019)]{education_lim2019} Genevieve Lim; A. Shelley; Dongcheol Heo. The regulation of learning and co-creation of new knowledge in mobile learning. Knowledge Management \& E-Learning: An International Journal, 2019. DOI: 10.34105/j.kmel.2019.11.024. URL: \url{https://doi.org/10.34105/j.kmel.2019.11.024}. Consulta: 5 out. 2026. Nível de acesso registrado na matriz de fontes.


\bibitem[Limone et al.(2020)]{education_limone2020} P. Limone; M. Sinatra; F. Ceglie; L. Monacis. Examining Procrastination among University Students through the Lens of the Self-Regulated Learning Model. Behavioral Sciences, 2020. DOI: 10.3390/bs10120184. URL: \url{https://doi.org/10.3390/bs10120184}. Consulta: 5 out. 2026. Nível de acesso registrado na matriz de fontes.


\bibitem[Palalas e Wark(2020)]{education_palalas2020} A. Palalas; Norine Wark. The relationship between mobile learning and self-regulated learning: A systematic review. Australasian Journal of Educational Technology, 2020. DOI: 10.14742/ajet.5650. URL: \url{https://doi.org/10.14742/ajet.5650}. Consulta: 5 out. 2026. Nível de acesso registrado na matriz de fontes.


\bibitem[Patzak et al.(2025)]{education_patzak2025} Alexandra Patzak; Xiaorong Zhang; Jovita Vytasek. Boosting productivity and wellbeing through time management: evidence-based strategies for higher education and workforce development. Frontiers in Education, 2025. DOI: 10.3389/feduc.2025.1623228. URL: \url{https://doi.org/10.3389/feduc.2025.1623228}. Consulta: 5 out. 2026. Nível de acesso registrado na matriz de fontes.


\bibitem[Prasse et al.(2024)]{education_prasse2024} Doreen Prasse; Mary Webb; Michelle Deschênes; Séverine Parent; Franziska Aeschlimann; Yoshiko Goda; Masanori Yamada; Audrey Raynault. Challenges in Promoting Self-Regulated Learning in Technology Supported Learning Environments: An Umbrella Review of Systematic Reviews and Meta-Analyses. Technology, Knowledge and Learning, 2024. DOI: 10.1007/s10758-024-09772-z. URL: \url{https://doi.org/10.1007/s10758-024-09772-z}. Consulta: 5 out. 2026. Nível de acesso registrado na matriz de fontes.


\bibitem[Taghavi-Nejad et al.(2024)]{education_taghavi2024} Fatemeh Sadat Taghavi-Nejad; Nasser Fallah; Behruz Lotfi Gaskaree. Mindfulness and Procrastination Among University EFL Learners: The Role of Attention Control and Self-Regulated Learning. Psychological Reports, 2024. DOI: 10.1177/00332941241287423. URL: \url{https://doi.org/10.1177/00332941241287423}. Consulta: 5 out. 2026. Nível de acesso registrado na matriz de fontes.


\bibitem[Tao et al.(2025)]{education_tao2025} Xue Tao; H. Hanif; Lie-Qin Wang. The effects of self-regulated learning strategies on academic procrastination and academic success among college EFL students in China. Frontiers in Psychology, 2025. DOI: 10.3389/fpsyg.2025.1562980. URL: \url{https://doi.org/10.3389/fpsyg.2025.1562980}. Consulta: 5 out. 2026. Nível de acesso registrado na matriz de fontes.


\bibitem[Wolters et al.(2017)]{education_wolters2017} C. Wolters; Sungjun Won; M. Hussain. Examining the relations of time management and procrastination within a model of self-regulated learning. Metacognition and Learning, 2017. DOI: 10.1007/s11409-017-9174-1. URL: \url{https://doi.org/10.1007/s11409-017-9174-1}. Consulta: 5 out. 2026. Nível de acesso registrado na matriz de fontes.


\bibitem[Wolters e Brady(2021)]{education_wolters2021} Christopher A. Wolters; Anna C. Brady. College Students’ Time Management: a Self-Regulated Learning Perspective. Educational Psychology Review, 2021. DOI: 10.1007/s10648-020-09519-z. URL: \url{https://doi.org/10.1007/s10648-020-09519-z}. Consulta: 5 out. 2026. Nível de acesso registrado na matriz de fontes.


\bibitem[Wong et al.(2026)]{education_wong2026} Jacqueline Wong; Mohammad Khalil; Vsevolod Suschevskiy; Martine Baars; Bjorn de Koning; Fred Paas. Student engagement profiles in a mobile app: Links to self-regulated learning and performance. Educational Technology Research and Development, 2026. DOI: 10.1007/s11423-026-10586-2. URL: \url{https://doi.org/10.1007/s11423-026-10586-2}. Consulta: 5 out. 2026. Nível de acesso registrado na matriz de fontes.


\bibitem[Xu et al.(2022)]{education_xu2022} Zhi-Hong Xu; Yingying Zhao; Bingsheng Zhang; J. Liew; Ashlynn Kogut. A meta-analysis of the efficacy of self-regulated learning interventions on academic achievement in online and blended environments in K-12 and higher education. Behaviour \& Information Technology, 2022. DOI: 10.1080/0144929X.2022.2151935. URL: \url{https://doi.org/10.1080/0144929X.2022.2151935}. Consulta: 5 out. 2026. Nível de acesso registrado na matriz de fontes.


\bibitem[Zheng et al.(2016)]{education_zheng2016} Lanqin Zheng; Xin Li; Fengying Chen. Effects of a mobile self-regulated learning approach on students’ learning achievements and self-regulated learning skills. Innovations in Education and Teaching International, 2016. DOI: 10.1080/14703297.2016.1259080. URL: \url{https://doi.org/10.1080/14703297.2016.1259080}. Consulta: 5 out. 2026. Nível de acesso registrado na matriz de fontes.


\bibitem[Azionya e Nhedzi(2021)]{azionya2021} C. Azionya; Abyshey Nhedzi. THE DIGITAL DIVIDE AND HIGHER EDUCATION CHALLENGE WITH EMERGENCY ONLINE LEARNING: ANALYSIS OF TWEETS IN THE WAKE OF THE COVID-19 LOCKDOWN. Turkish Online Journal of Distance Education, 2021. DOI: 10.17718/tojde.1002822. URL: \url{https://doi.org/10.17718/tojde.1002822}. Consulta: 5 out. 2026. Nível de acesso registrado na matriz de fontes.


\bibitem[Batanero-Ochaíta et al.(2021)]{batanero2021} Concepción Batanero-Ochaíta; Luis de-Marcos; Luis Felipe Rivera; J. Holvikivi; J. Hilera; Salvador Otón Tortosa. Improving Accessibility in Online Education: Comparative Analysis of Attitudes of Blind and Deaf Students Toward an Adapted Learning Platform. IEEE Access, 2021. DOI: 10.1109/ACCESS.2021.3095041. URL: \url{https://doi.org/10.1109/ACCESS.2021.3095041}. Consulta: 5 out. 2026. Nível de acesso registrado na matriz de fontes.


\bibitem[Bong e Chen(2021)]{bong2021} W. K. Bong; Wei-Qin Chen. Increasing faculty’s competence in digital accessibility for inclusive education: a systematic literature review. International Journal of Inclusive Education, 2021. DOI: 10.1080/13603116.2021.1937344. URL: \url{https://doi.org/10.1080/13603116.2021.1937344}. Consulta: 5 out. 2026. Nível de acesso registrado na matriz de fontes.


\bibitem[Choi e Seo(2024)]{choi2024} Gi Woong Choi; JooYoung Seo. Accessibility, Usability, and Universal Design for Learning: Discussion of Three Key LX/UX Elements for Inclusive Learning Design. TechTrends, 2024. DOI: 10.1007/s11528-024-00987-6. URL: \url{https://doi.org/10.1007/s11528-024-00987-6}. Consulta: 5 out. 2026. Nível de acesso registrado na matriz de fontes.


\bibitem[Cinquin et al.(2019)]{cinquin2019} Pierre-Antoine Cinquin; P. Guitton; H. Sauzéon. Online e-learning and cognitive disabilities: A systematic review. Comput. Educ., 2019. DOI: 10.1016/j.compedu.2018.12.004. URL: \url{https://doi.org/10.1016/j.compedu.2018.12.004}. Consulta: 5 out. 2026. Nível de acesso registrado na matriz de fontes.


\bibitem[Ahmad Faudzi et al.(2023)]{cob2023} M. Ahmad Faudzi; Zaihisma Che Cob; Ridha Omar; Sharul Azim Sharudin; Masitah Ghazali. Investigating the User Interface Design Frameworks of Current Mobile Learning Applications: A Systematic Review. Education Sciences, 2023. DOI: 10.3390/educsci13010094. URL: \url{https://doi.org/10.3390/educsci13010094}. Consulta: 5 out. 2026. Nível de acesso registrado na matriz de fontes.


\bibitem[Farhan e Razmak(2020)]{farhan2020} Wejdan Farhan; Jamil Razmak. A comparative study of an assistive e-learning interface among students with and without visual and hearing impairments. Disability and Rehabilitation: Assistive Technology, 2020. DOI: 10.1080/17483107.2020.1786733. URL: \url{https://doi.org/10.1080/17483107.2020.1786733}. Consulta: 5 out. 2026. Nível de acesso registrado na matriz de fontes.


\bibitem[Gobbo e Bezerra(2023)]{gobbo2023} José Alcides Gobbo; B. Bezerra. Emerging Themes for Digital Accessibility in Education. Sustainability, 2023. DOI: 10.3390/su151411392. URL: \url{https://doi.org/10.3390/su151411392}. Consulta: 5 out. 2026. Nível de acesso registrado na matriz de fontes.


\bibitem[Guo e Wan(2022)]{guo2022} Congbin Guo; Boshen Wan. The digital divide in online learning in China during the COVID-19 pandemic. Technology in Society, 2022. DOI: 10.1016/j.techsoc.2022.102122. URL: \url{https://doi.org/10.1016/j.techsoc.2022.102122}. Consulta: 5 out. 2026. Nível de acesso registrado na matriz de fontes.


\bibitem[Jamil e Muschert(2023)]{jamil2023} S. Jamil; G. Muschert. The COVID-19 Pandemic and E-Learning: The Digital Divide and Educational Crises in Pakistan’s Universities. The American Behavioral Scientist, 2023. DOI: 10.1177/00027642231156779. URL: \url{https://doi.org/10.1177/00027642231156779}. Consulta: 5 out. 2026. Nível de acesso registrado na matriz de fontes.


\bibitem[Kumar e Mohite(2017)]{kumar2017} B. Kumar; P. Mohite. Usability of mobile learning applications: a systematic literature review. Journal of Computers in Education, 2017. DOI: 10.1007/s40692-017-0093-6. URL: \url{https://doi.org/10.1007/s40692-017-0093-6}. Consulta: 5 out. 2026. Nível de acesso registrado na matriz de fontes.


\bibitem[Kumar et al.(2019)]{kumarguideline2019} B. Kumar; M. S. Goundar; S. Chand. Usability guideline for Mobile learning applications: an update. Education and Information Technologies, 2019. DOI: 10.1007/s10639-019-09937-9. URL: \url{https://doi.org/10.1007/s10639-019-09937-9}. Consulta: 5 out. 2026. Nível de acesso registrado na matriz de fontes.


\bibitem[Kumar et al.(2024)]{kumarmapping2024} B. Kumar; S. Chand; M. S. Goundar. Usability testing of mobile learning applications: a systematic mapping study. The International Journal of Information and Learning Technology, 2024. DOI: 10.1108/IJILT-03-2023-0029. URL: \url{https://doi.org/10.1108/IJILT-03-2023-0029}. Consulta: 5 out. 2026. Nível de acesso registrado na matriz de fontes.


\bibitem[Kumi-Yeboah et al.(2023)]{kumiyeboah2023} Alex Kumi-Yeboah; YangHyun Kim; Yaa Essah Armah. Strategies for overcoming the digital divide during the COVID-19 pandemic in higher education institutions in Ghana. Br. J. Educ. Technol., 2023. DOI: 10.1111/bjet.13356. URL: \url{https://doi.org/10.1111/bjet.13356}. Consulta: 5 out. 2026. Nível de acesso registrado na matriz de fontes.


\bibitem[Linhalis et al.(2026)]{linhalis2026} Flávia Linhalis; Thiago Máximo Pavão; André Constantino da Silva. Educa Offline: concepción de un Ambiente Virtual de Aprendizaje (AVA) en contextos de baja conectividad y vulnerabilidad social. Revista Latinoamericana de Tecnología Educativa - RELATEC, 2026. DOI: 10.17398/1695-288X.25.1.59. URL: \url{https://doi.org/10.17398/1695-288X.25.1.59}. Consulta: 5 out. 2026. Nível de acesso registrado na matriz de fontes.


\bibitem[Martin et al.(2024)]{martin2024} Florence Martin; E. Ceviker; T. Gezer. From digital divide to digital equity: Systematic review of two decades of research on educational digital divide factors, dimensions, and interventions. Journal of Research on Technology in Education, 2024. DOI: 10.1080/15391523.2024.2425442. URL: \url{https://doi.org/10.1080/15391523.2024.2425442}. Consulta: 5 out. 2026. Nível de acesso registrado na matriz de fontes.


\bibitem[Mathrani et al.(2021)]{mathrani2021} A. Mathrani; Tarushikha Sarvesh; R. Umer. Digital divide framework: online learning in developing countries during the COVID-19 lockdown. Globalisation, Societies and Education, 2021. DOI: 10.1080/14767724.2021.1981253. URL: \url{https://doi.org/10.1080/14767724.2021.1981253}. Consulta: 5 out. 2026. Nível de acesso registrado na matriz de fontes.


\bibitem[Naveed et al.(2023)]{naveed2023} Q. Naveed; Heena Choudhary; Naim Ahmad; Jarallah Alqahtani; A. Qahmash. Mobile Learning in Higher Education: A Systematic Literature Review. Sustainability, 2023. DOI: 10.3390/su151813566. URL: \url{https://doi.org/10.3390/su151813566}. Consulta: 5 out. 2026. Nível de acesso registrado na matriz de fontes.


\bibitem[Punchoojit e Hongwarittorrn(2017)]{punchoojit2017} Lumpapun Punchoojit; Nuttanont Hongwarittorrn. Usability Studies on Mobile User Interface Design Patterns: A Systematic Literature Review. Advances in Human-Computer Interaction, 2017. DOI: 10.1155/2017/6787504. URL: \url{https://doi.org/10.1155/2017/6787504}. Consulta: 5 out. 2026. Nível de acesso registrado na matriz de fontes.


\bibitem[Werfhorst et al.(2022)]{werfhorst2022} H. G. van de Werfhorst; Emma Kessenich; Sara Geven. The digital divide in online education: Inequality in digital readiness of students and schools. Computers and Education Open, 2022. DOI: 10.1016/j.caeo.2022.100100. URL: \url{https://doi.org/10.1016/j.caeo.2022.100100}. Consulta: 5 out. 2026. Nível de acesso registrado na matriz de fontes.


\bibitem[Squarcina et al.(2021)]{OS01} Marco Squarcina; Stefano Calzavara; Matteo Maffei. The Remote on the Local: Exacerbating Web Attacks Via Service Workers Caches. 2021 IEEE Security and Privacy Workshops (SPW), 2021. DOI: 10.1109/SPW53761.2021.00062. URL: \url{https://doi.org/10.1109/SPW53761.2021.00062}. Consulta: 5 out. 2026. Nível de acesso registrado na matriz de fontes.


\bibitem[Subramani et al.(2022)]{OS02} Karthika Subramani; Jordan Jueckstock; Alexandros Kapravelos; Roberto Perdisci. SoK: Workerounds - Categorizing Service Worker Attacks and Mitigations. 2022 IEEE 7th European Symposium on Security and Privacy (EuroS\&P), 2022. DOI: 10.1109/EuroSP53844.2022.00041. URL: \url{https://doi.org/10.1109/EuroSP53844.2022.00041}. Consulta: 5 out. 2026. Nível de acesso registrado na matriz de fontes.


\bibitem[Karami et al.(2021)]{OS03} Soroush Karami; Panagiotis Ilia; Jason Polakis. Awakening the Web's Sleeper Agents: Misusing Service Workers for Privacy Leakage.. Network and Distributed System Security Symposium (NDSS), 2021, 2021. DOI: 10.14722/NDSS.2021.23104. URL: \url{https://doi.org/10.14722/NDSS.2021.23104}. Consulta: 5 out. 2026. Nível de acesso registrado na matriz de fontes.


\bibitem[Dehling et al.(2019)]{OS06} Florian Dehling; Tobias Mengel; Luigi Lo Iacono. Rotten Cellar: Security and Privacy of the Browser Cache Revisited. Secure IT Systems: 24th Nordic Conference, NordSec 2019, pp.20–36, 2019. DOI: 10.1007/978-3-030-35055-0\_2. URL: \url{https://doi.org/10.1007/978-3-030-35055-0_2}. Consulta: 5 out. 2026. Nível de acesso registrado na matriz de fontes.


\bibitem[Haas et al.(2023)]{OS07} Julian Haas; Ragnar Mogk; Elena Yanakieva; Annette Bieniusa; Mira Mezini. LoRe: A Programming Model for Verifiably Safe Local-First Software. ACM Transactions on Programming Languages and Systems, 2023. DOI: 10.1145/3633769. URL: \url{https://doi.org/10.1145/3633769}. Consulta: 5 out. 2026. Nível de acesso registrado na matriz de fontes.


\bibitem[Litt e Jackson(2022)]{OS08} Geoffrey Litt; B. Jackson. Merge what you can, fork what you can't: managing data integrity in local-first software. 9th Workshop on Principles and Practice of Consistency for Distributed Data (PaPoC), 2022, 2022. DOI: 10.1145/3517209.3524041. URL: \url{https://doi.org/10.1145/3517209.3524041}. Consulta: 5 out. 2026. Nível de acesso registrado na matriz de fontes.


\bibitem[MacFadden e Qiu(2022)]{OS09} Michael MacFadden; Meikang Qiu. Performance Impacts of JavaScript-Based Encryption of HTML5 Web Storage for Enhanced Privacy. IEEE SmartCloud 2022, 2022. DOI: 10.1109/SmartCloud55982.2022.00037. URL: \url{https://doi.org/10.1109/SmartCloud55982.2022.00037}. Consulta: 5 out. 2026. Nível de acesso registrado na matriz de fontes.


\bibitem[Jannes et al.(2023)]{OS10} Kristof Jannes; Emad Heydari Beni; Bert Lagaisse; Wouter Joosen. BeauForT: Robust Byzantine Fault Tolerance for Client-Centric Mobile Web Applications. IEEE Transactions on Parallel and Distributed Systems, 2023. DOI: 10.1109/TPDS.2023.3241963. URL: \url{https://doi.org/10.1109/TPDS.2023.3241963}. Consulta: 5 out. 2026. Nível de acesso registrado na matriz de fontes.


\bibitem[Ekambaranathan et al.(2023)]{OS13} Anirudh Ekambaranathan; Jun Zhao; George Chalhoub. Navigating the Data Avalanche: Towards Supporting Developers in Developing Privacy-Friendly Children’s Apps. Proceedings of the ACM on interactive, mobile, wearable and ubiquitous technologies, 2023. DOI: 10.1145/3596267. URL: \url{https://doi.org/10.1145/3596267}. Consulta: 5 out. 2026. Nível de acesso registrado na matriz de fontes.


\bibitem[Gruber et al.(2022)]{OS14} Moritz Gruber; Christian Höfig; Maximilian Golla; Tobias Urban; Matteo Große-Kampmann. "We may share the number of diaper changes": A Privacy and Security Analysis of Mobile Child Care Applications. Proceedings on Privacy Enhancing Technologies, 2022. DOI: 10.56553/popets-2022-0078. URL: \url{https://doi.org/10.56553/popets-2022-0078}. Consulta: 5 out. 2026. Nível de acesso registrado na matriz de fontes.


\bibitem[Zhao et al.(2022)]{OS15} Jun Zhao; Blanche Duron; Ge Wang. KOALA Hero: Inform Children of Privacy Risks of Mobile Apps. Interaction Design and Children, 2022, 2022. DOI: 10.1145/3501712.3535278. URL: \url{https://doi.org/10.1145/3501712.3535278}. Consulta: 5 out. 2026. Nível de acesso registrado na matriz de fontes.


\bibitem[Sun et al.(2023)]{OS17} Ruoxi Sun; Minhui Xue; Gareth Tyson; Shuang Wang; Seyit Camtepe; Surya Nepal. Not Seen, Not Heard in the Digital World! Measuring Privacy Practices in Children’s Apps. ACM Web Conference 2023, pp.2166–2177, 2023. DOI: 10.1145/3543507.3583327. URL: \url{https://doi.org/10.1145/3543507.3583327}. Consulta: 5 out. 2026. Nível de acesso registrado na matriz de fontes.


\bibitem[Mkpojiogu et al.(2021)]{OS18} Emmanuel O.C. Mkpojiogu; Azham Hussain; Monday Onah Agbudu. Security Issues in the Use of Mobile Educational Apps: A Review. International Journal of Interactive Mobile Technologies (ijim), 2021. DOI: 10.3991/IJIM.V15I06.20631. URL: \url{https://doi.org/10.3991/IJIM.V15I06.20631}. Consulta: 5 out. 2026. Nível de acesso registrado na matriz de fontes.


\bibitem[Kong et al.(2019)]{eng01} Pingfan Kong; Li Li; Jun Gao; Kui Liu; Tegawendé F. Bissyandé; Jacques Klein. Automated Testing of Android Apps: A Systematic Literature Review. IEEE Transactions on Reliability, 2019. DOI: 10.1109/TR.2018.2865733. URL: \url{https://doi.org/10.1109/TR.2018.2865733}. Consulta: 5 out. 2026. Nível de acesso registrado na matriz de fontes.


\bibitem[Berihun et al.(2023)]{eng02} Natnael Gonfa Berihun; Cyrille Dongmo; J. A. van der Poll. The Applicability of Automated Testing Frameworks for Mobile Application Testing: A Systematic Literature Review. Computers, 2023. DOI: 10.3390/computers12050097. URL: \url{https://doi.org/10.3390/computers12050097}. Consulta: 5 out. 2026. Nível de acesso registrado na matriz de fontes.


\bibitem[Tramontana et al.(2018)]{eng03} Porfirio Tramontana; Domenico Amalfitano; Nicola Amatucci; A. Fasolino. Automated functional testing of mobile applications: a systematic mapping study. Software Quality Journal, 2018. DOI: 10.1007/s11219-018-9418-6. URL: \url{https://doi.org/10.1007/s11219-018-9418-6}. Consulta: 5 out. 2026. Nível de acesso registrado na matriz de fontes.


\bibitem[Zein et al.(2023)]{eng04} Samer Zein; Norsaremah Salleh; John C. Grundy. Systematic reviews in mobile app software engineering: A tertiary study. Inf. Softw. Technol., 2023. DOI: 10.1016/j.infsof.2023.107323. URL: \url{https://doi.org/10.1016/j.infsof.2023.107323}. Consulta: 5 out. 2026. Nível de acesso registrado na matriz de fontes.


\bibitem[Zein et al.(2016)]{eng05} Samer Zein; Norsaremah Salleh; J. Grundy. A systematic mapping study of mobile application testing techniques. J. Syst. Softw., 2016. DOI: 10.1016/j.jss.2016.03.065. URL: \url{https://doi.org/10.1016/j.jss.2016.03.065}. Consulta: 5 out. 2026. Nível de acesso registrado na matriz de fontes.


\bibitem[Júnior et al.(2022)]{eng06} Misael C. Júnior; Domenico Amalfitano; Lina Garcés; A. Fasolino; Stevão A. Andrade; M. Delamaro. Dynamic Testing Techniques of Non-functional Requirements in Mobile Apps: A Systematic Mapping Study. ACM Computing Surveys (CSUR), 2022. DOI: 10.1145/3507903. URL: \url{https://doi.org/10.1145/3507903}. Consulta: 5 out. 2026. Nível de acesso registrado na matriz de fontes.


\bibitem[Weigand e Johannesson(2023)]{eng07} H. Weigand; P. Johannesson. How to Identify your Design Science Research Artifact. 2023 IEEE 25th Conference on Business Informatics (CBI), 2023. DOI: 10.1109/CBI58679.2023.10187511. URL: \url{https://doi.org/10.1109/CBI58679.2023.10187511}. Consulta: 5 out. 2026. Nível de acesso registrado na matriz de fontes.


\bibitem[Gonzalez-Barahona e Robles(2023)]{eng08} Jesus M. Gonzalez-Barahona; G. Robles. Revisiting the reproducibility of empirical software engineering studies based on data retrieved from development repositories. Inf. Softw. Technol., 2023. DOI: 10.1016/j.infsof.2023.107318. URL: \url{https://doi.org/10.1016/j.infsof.2023.107318}. Consulta: 5 out. 2026. Nível de acesso registrado na matriz de fontes.


\bibitem[Liu et al.(2022)]{eng10} Chao Liu; Cui-Yun Gao; Xin Xia; David Lo; John C. Grundy; Xiao-Hu Yang. On the Reproducibility and Replicability of Deep Learning in Software Engineering. ACM Trans. Softw. Eng. Methodol., 2022. DOI: 10.1145/3477535. URL: \url{https://doi.org/10.1145/3477535}. Consulta: 5 out. 2026. Nível de acesso registrado na matriz de fontes.


\bibitem[Anchundia e Fonseca C.(2020)]{eng11} Carlos E. Anchundia; Efraín R. Fonseca C.. Resources for Reproducibility of Experiments in Empirical Software Engineering: Topics Derived From a Secondary Study. IEEE Access, 2020. DOI: 10.1109/ACCESS.2020.2964587. URL: \url{https://doi.org/10.1109/ACCESS.2020.2964587}. Consulta: 5 out. 2026. Nível de acesso registrado na matriz de fontes.


\bibitem[Santos et al.(2022)]{eng12} Daniel Amador dos Santos; E. D. de Almeida; Iftekhar Ahmed. Investigating replication challenges through multiple replications of an experiment. Inf. Softw. Technol., 2022. DOI: 10.1016/j.infsof.2022.106870. URL: \url{https://doi.org/10.1016/j.infsof.2022.106870}. Consulta: 5 out. 2026. Nível de acesso registrado na matriz de fontes.


\bibitem[Shepperd et al.(2018)]{eng13} M. Shepperd; N. Ajienka; S. Counsell. The role and value of replication in empirical software engineering results. Inf. Softw. Technol., 2018. DOI: 10.1016/j.infsof.2018.01.006. URL: \url{https://doi.org/10.1016/j.infsof.2018.01.006}. Consulta: 5 out. 2026. Nível de acesso registrado na matriz de fontes.


\bibitem[Shahin et al.(2017)]{eng14} Mojtaba Shahin; Muhammad Ali Babar; Li-Ming Zhu. Continuous Integration, Delivery and Deployment: A Systematic Review on Approaches, Tools, Challenges and Practices. IEEE Access, 2017. DOI: 10.1109/ACCESS.2017.2685629. URL: \url{https://doi.org/10.1109/ACCESS.2017.2685629}. Consulta: 5 out. 2026. Nível de acesso registrado na matriz de fontes.


\bibitem[Soares et al.(2022)]{eng15} Eliezio Soares; Gustavo Sizílio; Jadson Santos; Daniel Alencar da Costa; U. Kulesza. The effects of continuous integration on software development: a systematic literature review. Empirical Software Engineering, 2022. DOI: 10.1007/s10664-021-10114-1. URL: \url{https://doi.org/10.1007/s10664-021-10114-1}. Consulta: 5 out. 2026. Nível de acesso registrado na matriz de fontes.


\bibitem[Elazhary et al.(2021)]{eng16} Omar Elazhary; Colin M. Werner; Ze-Shi Li; Derek Lowlind; Neil A. Ernst; M. Storey. Uncovering the Benefits and Challenges of Continuous Integration Practices. IEEE Transactions on Software Engineering, 2021. DOI: 10.1109/TSE.2021.3064953. URL: \url{https://doi.org/10.1109/TSE.2021.3064953}. Consulta: 5 out. 2026. Nível de acesso registrado na matriz de fontes.


\bibitem[Guevara-Vega et al.(2024)]{eng18} C. Guevara-Vega; B. Bernárdez; Margarita Cruz; A. Durán; Antonio Ruiz-Cortés; Martin Solari. Research artifacts for human-oriented experiments in software engineering: An ACM badges-driven structure proposal. J. Syst. Softw., 2024. DOI: 10.1016/j.jss.2024.112187. URL: \url{https://doi.org/10.1016/j.jss.2024.112187}. Consulta: 5 out. 2026. Nível de acesso registrado na matriz de fontes.


\bibitem[Baskerville et al.(2026)]{eng20} Richard Baskerville; Jan Pries-Heje; John R. Venable. MEDS: Methodology for Evaluation in Design Science. European Journal of Information Systems, 2026. DOI: 10.1080/0960085X.2026.2627280. URL: \url{https://doi.org/10.1080/0960085X.2026.2627280}. Consulta: 5 out. 2026. Nível de acesso registrado na matriz de fontes.


\bibitem[Page et al.(2021a)]{SU01} Matthew J. Page; Joanne E. McKenzie; Patrick M.M. Bossuyt; Isabelle Boutron; Tammy Hoffmann; Cynthia D. Mulrow; Larissa Shamseer; Jennifer Tetzlaff; Elie A. Akl; Sue E. Brennan; Roger Chou; Julie Glanville; Jeremy M. Grimshaw; Asbjørn Hróbjartsson; Manoj M. Lalu; Tianjing Li; Elizabeth Loder; Evan Mayo-Wilson; Steve McDonald; Luke A McGuinness; Lesley A. Stewart; James Thomas; Andrea C. Tricco; Vivian Welch; Penny Whiting; David Moher. The PRISMA 2020 statement: an updated guideline for reporting systematic reviews. BMJ, 2021. DOI: 10.1136/bmj.n71. URL: \url{https://doi.org/10.1136/bmj.n71}. Consulta: 5 out. 2026. Nível de acesso registrado na matriz de fontes.


\bibitem[Page et al.(2021b)]{SU02} Matthew J. Page; David Moher; Patrick M.M. Bossuyt; Isabelle Boutron; Tammy Hoffmann; Cynthia D. Mulrow; Larissa Shamseer; Jennifer Tetzlaff; Elie A. Akl; Sue E. Brennan; Roger Chou; Julie Glanville; Jeremy M. Grimshaw; Asbjørn Hróbjartsson; Manoj M. Lalu; Tianjing Li; Elizabeth Loder; Evan Mayo-Wilson; Steve McDonald; Luke A McGuinness; Lesley A. Stewart; James Thomas; Andrea C. Tricco; Vivian Welch; Penny Whiting; Joanne E. McKenzie. PRISMA 2020 explanation and elaboration: updated guidance and exemplars for reporting systematic reviews.. BMJ, 2021. DOI: 10.1136/bmj.n160. URL: \url{https://doi.org/10.1136/bmj.n160}. Consulta: 5 out. 2026. Nível de acesso registrado na matriz de fontes.


\bibitem[Rethlefsen et al.(2021)]{SU03} Melissa L. Rethlefsen; Shona Kirtley; Siw Waffenschmidt; Ana Patricia Ayala; David Moher; Matthew J. Page; Jonathan Koffel; PRISMA-S Group. PRISMA-S: an extension to the PRISMA Statement for Reporting Literature Searches in Systematic Reviews.. Systematic Reviews, 2021. DOI: 10.1186/s13643-020-01542-z. URL: \url{https://doi.org/10.1186/s13643-020-01542-z}. Consulta: 5 out. 2026. Nível de acesso registrado na matriz de fontes.


\bibitem[Munafò et al.(2017)]{SU04} Marcus R. Munafò; Brian A. Nosek; Dorothy V. M. Bishop; Katherine S. Button; Christopher D. Chambers; Nathalie Percie du Sert; Uri Simonsohn; Eric-Jan Wagenmakers; Jennifer J. Ware; John P. A. Ioannidis. A manifesto for reproducible science. Nature Human Behaviour, 2017. DOI: 10.1038/s41562-016-0021. URL: \url{https://doi.org/10.1038/s41562-016-0021}. Consulta: 5 out. 2026. Nível de acesso registrado na matriz de fontes.


\bibitem[Foltýnek et al.(2019)]{SU05} Tomáš Foltýnek; Norman Meuschke; Bela Gipp. Academic Plagiarism Detection: A Systematic Literature Review. ACM Computing Surveys, 2019. DOI: 10.1145/3345317. URL: \url{https://doi.org/10.1145/3345317}. Consulta: 5 out. 2026. Nível de acesso registrado na matriz de fontes.


\bibitem[Walters e Wilder(2023)]{SU06} William H. Walters; Esther Isabelle Wilder. Fabrication and errors in the bibliographic citations generated by ChatGPT. Scientific Reports, 13, article 14045, 2023. DOI: 10.1038/s41598-023-41032-5. URL: \url{https://doi.org/10.1038/s41598-023-41032-5}. Consulta: 5 out. 2026. Nível de acesso registrado na matriz de fontes.


\bibitem[Bhattacharyya et al.(2023)]{SU07} Mehul Bhattacharyya; Valerie M. Miller; Debjani Bhattacharyya; Larry E. Miller. High Rates of Fabricated and Inaccurate References in ChatGPT-Generated Medical Content. Cureus, 2023. DOI: 10.7759/cureus.39238. URL: \url{https://doi.org/10.7759/cureus.39238}. Consulta: 5 out. 2026. Nível de acesso registrado na matriz de fontes.


\bibitem[Dipongkor et al.(2021)]{SU08} Atish Kumar Dipongkor; Rayhanul Islam; Md. Shafiuzzaman; Asif Nashiry; Syed Md. Galib; Khaza Moinuddin Mazumder. AcPgChecker: Detection of Plagiarism among Academic and Scientific Writings. 2021 Joint 10th International Conference on Informatics, Electronics \& Vision (ICIEV) and 2021 5th International Conference on Imaging, Vision \& Pattern Recognition (icIVPR), 2021. DOI: 10.1109/icievicivpr52578.2021.9564123. URL: \url{https://doi.org/10.1109/icievicivpr52578.2021.9564123}. Consulta: 5 out. 2026. Nível de acesso registrado na matriz de fontes.


\bibitem[Panadero(2017)]{SU10} Ernesto Panadero. A Review of Self-regulated Learning: Six Models and Four Directions for Research. Frontiers in Psychology, 2017. DOI: 10.3389/fpsyg.2017.00422. URL: \url{https://doi.org/10.3389/fpsyg.2017.00422}. Consulta: 5 out. 2026. Nível de acesso registrado na matriz de fontes.


\bibitem[Snyder(2019)]{SU11} Hannah Snyder. Literature review as a research methodology: An overview and guidelines. Journal of Business Research, 2019. DOI: 10.1016/j.jbusres.2019.07.039. URL: \url{https://doi.org/10.1016/j.jbusres.2019.07.039}. Consulta: 5 out. 2026. Nível de acesso registrado na matriz de fontes.

\end{thebibliography}
